\section{Testdokumentation}
Tests müssen aus verschiedenen Gründen \textbf{dokumentiert} werden:

\begin{itemize}
    \item Tests sollen von Mitarbeitern und Kunden begutachtet werden können, damit diese feststellen können, ob Annahmen der Tester korrekt und Tests vollständig sind; Projektleiter möchten Testfälle prüfen, damit Merkmale mit hoher Priorität genügend intensiv und Merkmale mit niedriger Priorität nicht zu intensiv geprüft werden
    \item die Dokumentation soll helfen, den Testbetrieb reibungslos ablaufen zu lassen, für den meist wenig Zeit zur Verfügung steht
    \item die Dokumentation soll Testern helfen, wenn eine Unterbrechung der Durchführung der Tests erfolgen muss (aufgrund Fehler oder Ausfall Personal)
    \item ein schriftlicher Nachweis der Durchführung von Tests ist oft Voraussetzung für eine \textbf{Zertifizierung}, bspw. nach \textbf{ISO 9000}\footnote{
    \url{https://de.m.wikipedia.org/wiki/Qualit\%C3\%A4tsmanagementnorm}, abgerufen 01.06.2024
    } oder \textbf{CMM}\footnote{
    \url{https://de.wikipedia.org/wiki/Capability_Maturity_Model}, abgerufen 01.06.2024
    }
\end{itemize}

\noindent
In der Praxis hat sich eine Unterteilung der Dokumentation in \textbf{Testfälle}, \textbf{Testprotokolle} und \textbf{Fehlerberichte} als geeignet erwiesen:

\begin{itemize}
    \item der \textbf{Testfall} beschreibt die Durchführung eines Tests detailliert
    \item das \textbf{Testprotokoll} protokolliert die Durchführung eines Tests
    \item im Test gefundene Fehler werden in \textbf{Fehlerberichten} detailliert festgehalten und verarbeitet
\end{itemize}

\noindent
Dieses Vorgehen hat verschiedene Gründe: Die Aufteilung in \textbf{Testfall} und \textbf{Testprotokoll} ist wichtig, da ein Test meist mehrfach durchgeführt wird, bspw. wenn Fehler auftreten oder mehrere Lieferungen erfolgen - der gleiche Test muss also mehrfach dokumentiert werden, was mit \textit{minimalem Aufwand} gelingt, wenn Testfall und Protokoll getrennt sind.\\
Die gesonderte Verarbeitung von \textbf{Fehlerberichten} ist günstig, da einerseits ein Fehler gleichzeitig in mehreren Tests auftreten kann (s. Abbildung~\ref{fig:testdokumentation}) und andererseits Fehlerberichte Gegenstand eines komplexen Arbeitsprozesses sind.

\begin{figure}
    \centering
    \includegraphics[scale=0.4]{part four/Management des QS-Prozesses/img/testdokumentation}
    \caption{Übersicht über die Testdokumentation. (Quelle: in Anlehnung an~\cite[Abb. 6.1, 72]{Wed09c})}
    \label{fig:testdokumentation}
\end{figure}


\subsubsection*{Testfall}
Ein \textbf{Testfall} beschreibt detailliert die Durchführung von Tests und dient vor allem dazu, dass Tester die Tests schnell und ohne besonderen Vorkenntnisse der fachlichen Zusammenhänge durchführen können.\\
Hierzu werden die durchzuführenden Aktionen und die erwarteten Ergebnisse tabellarisch beschrieben, wie im folgenden Beispiel:

\subsubsection*{Testfall MAR1-01: Kundenliste erzeugen für PLZ - Gutfall}

    \begin{itemize}
        \item[] \textbf{Autor}: at
        \item[] \textbf{Kategorie}: Systemtest
        \item[] \textbf{Beschreibung}:
        \item[] \textbf{Setup}:
        \begin{itemize}
            \item Im System sind die Kundendaten TD01 eingespielt und nicht modifiziert worden
            \item Der Nutzer ist angemeldet und besitzt die Rolle ``Leitender Mitarbeiter``
        \end{itemize}
    \end{itemize}

    \noindent
    \textbf{Vorgehen}:
    \begin{table}[]
        \centering
        \setlength{\tabcolsep}{0.5em}
        \def\arraystretch{1.5}
        \begin{tabular}{|c|l|c|}
            \hline
            \multicolumn{1}{|l|}{\textbf{Lfd. Nr.}} & \multicolumn{1}{c|}{\textbf{Schritt}}        & \textbf{erwartetes Ergebnis}                        \\ \hline
            1                                       & Auswahl PLZ Combobox                         & ist ausgewählt                                      \\ \hline
            2                                       & Auswahl von-bis                              & ist ausgewählt                                      \\ \hline
            3                                       & Eingabe 18430 und 18435 in die Eingabefelder & Zahlen werden dargestellt                           \\ \hline
            4                              & Export anklicken                             & \multicolumn{1}{l|}{Dateiauswahldialog öffnet sich} \\ \hline
            5                              & Auswahl eines geeigneten Dateinamens         & \multicolumn{1}{l|}{Datei mit Inhalt der Daten}     \\ \hline
        \end{tabular}
    \end{table}

\begin{itemize}
    \item[] \textbf{Zugrundeliegende Anforderungen}:
    \begin{itemize}
        \item Anwendungsfall MAR-1 vom 11.01.08
        \item GUI-Entwurf MAR-1 vom 26.03.08
        \item Liste der Geschäftsregeln (GSR) vom 18.01.08
    \end{itemize}
\end{itemize}

\subsubsection*{Testprotokoll}
Im \textbf{Tesprotokoll} wird der \textbf{Testfortschritt} protokolliert - da der Ablauf bereits im \textbf{Testfall} beschrieben ist, kann das Protokoll kurz gefasst sein, wie im Folgenden Beispiel:

\subsubsection*{Testprotokoll für Testfall MAR1-01: Kundenliste erzeugen für PLZ - Gutfall}
\begin{table}[]
    \begin{tabular}{|c|c|c|c|c|c|c|}
        \hline
        \multicolumn{1}{|l|}{\textbf{Lfd. Nr.}} & \textbf{getestete Version} & \textbf{Datum} & \multicolumn{1}{l|}{\textbf{Tester}} & \multicolumn{1}{l|}{\textbf{bis Schritt}} & \multicolumn{1}{l|}{\textbf{Status}} & \multicolumn{1}{l|}{\textbf{Ticket}} \\ \hline
        1                                       & SYSTST2                    & 04.04.08       & at                                   & 1                                         & Fehler                               & 123                                  \\ \hline
        2                                       & SYSTST3                    & 04.04.08       & at                                   & 5                                         & Fehler                               & 456                                  \\ \hline
        3                                       & SYSTST5                    & 04.04.08       & at                                   & 5                                         & OK                                   &                                      \\ \hline
    \end{tabular}
\end{table}

\subsubsection*{Modultest}
Die Dokumentation von \textbf{Modultest} stellt einen Sonderfall dar.\\
Bei der Verwendung von Unittests (bspw. \textbf{JUnit}) kann die Dokumentation in dem Quellcode als Javadoc hinterlegt werden; außerdem ist eine Protokollierung nicht notwendig, da die Tests jederzeit ohne großen Aufwand durchgeführt werden können.